\maketitle
\makesignature

\ifproject
\begin{abstract}

โครงงานนี้นำเสนอการพัฒนาเกมต้นแบบแนวสองมิติประเภท Action RPG ผสมผสานกับรูปแบบการเล่นแบบ Casual ภายใต้ชื่อเกม Marionette ซึ่งมีเนื้อเรื่องเกี่ยวกับตัวละครเอกที่ถูกแม่มดจับตัวและสาปให้กลายเป็นหุ่นเชิดในโรงละครลึกลับ ผู้เล่นจะต้องพยายามหาทางหลบหนีจากการควบคุมของแม่มด พร้อมทั้งเผชิญหน้ากับศัตรูและอุปสรรคต่าง ๆ ภายในเกม รูปแบบการเล่นเน้นการวางแผนเชิงกลยุทธ์และการบริหารทรัพยากรที่มีอยู่อย่างจำกัดภายใต้สถานการณ์ที่มีข้อจำกัดด้านเวลา เพื่อให้สามารถรับมือกับศัตรูและผ่านด่านต่าง ๆ ได้สำเร็จ

จากรายงาน ดัชนีอุตสาหกรรมดิจิทัลประเทศไทย ปี 2567 ของสำนักงานส่งเสริมเศรษฐกิจดิจิทัล depa Thailand พบว่าแม้อุตสาหกรรมเกมของประเทศไทยจะมีมูลค่าเพิ่มขึ้นอย่างต่อเนื่อง แต่การพัฒนาเกมภายใต้ทรัพย์สินทางปัญญาของผู้พัฒนาในประเทศกลับมีแนวโน้มลดลง โดยส่วนหนึ่งมาจากการรับจ้างพัฒนาและจัดจำหน่ายเกมภายใต้ลิขสิทธิ์ของต่างประเทศ โครงงานนี้จึงมีเป้าหมายเพื่อส่งเสริมการพัฒนาเกมที่เป็นทรัพย์สินทางปัญญาของนักพัฒนาไทย และสร้างแรงบันดาลใจให้เกิดการพัฒนาเกมต้นฉบับภายในประเทศมากยิ่งขึ้น

การพัฒนาเกมต้นแบบนี้ใช้ Unity Game Engine เป็นเครื่องมือหลัก และใช้ภาษา C\# สำหรับการพัฒนาระบบต่าง ๆ ของเกม โดยมีการใช้ Animation State Machine เพื่อควบคุมการทำงานของแอนิเมชันตัวละคร และ Finite State Machine สำหรับกำหนดพฤติกรรมของตัวละครที่ไม่ใช่ผู้เล่น นอกจากนี้ยังมีการออกแบบส่วนติดต่อผู้ใช้ เอฟเฟกต์ภาพ เสียงประกอบ และดนตรีพื้นหลัง เพื่อเพิ่มประสบการณ์การเล่นเกมให้สมบูรณ์มากยิ่งขึ้น

\end{abstract}

\begin{abstract}
This project presents the development of a two-dimensional Action RPG combined with casual gameplay elements, titled Marionette. The game follows the story of a protagonist who is captured by a witch and cursed to become a marionette puppet in a mysterious theater. The player must attempt to escape from the witch’s control while facing various enemies and challenges throughout the game. The gameplay emphasizes strategic thinking and resource management within time-limited situations in order to overcome obstacles and survive encounters with enemies.

According to the report Thailand Digital Industry Index 2024 published by the Digital Economy Promotion Agency depa Thailand, the value of the Thai game industry has continued to grow. However, the development of original intellectual property (IP) by Thai companies has declined, with many companies focusing on publishing or outsourcing under foreign intellectual property licenses. This project therefore aims to encourage the development of original game content and inspire Thai developers to create games under their own intellectual property.

The game prototype was developed using the Unity Game Engine and the C\# programming language. Key development techniques include the use of Animation State Machines for character animations and Finite State Machines for non-player character behavior. Visual effects, sound effects, and background music are also integrated to enhance the gameplay experience. The resulting prototype demonstrates the feasibility of combining exploration, combat, and mini-game mechanics within a theatrical game theme.
\end{abstract}

\iffalse
\begin{dedication}
This document is dedicated to all Chiang Mai University students.

Dedication page is optional.
\end{dedication}
\fi % \iffalse

\iffalse
\begin{acknowledgments}
Your acknowledgments go here. Make sure it sits inside the
\texttt{acknowledgment} environment.

\acksign{2020}{5}{25}
\end{acknowledgments}%
\fi

\contentspage

\ifproject
%%\figurelistpage

%%\tablelistpage
%%\fi % \ifproject

% \abbrlist % this page is optional

% \symlist % this page is optional

% \preface % this section is optional
